\begin{description}
\item[Convention:] During this Section, we will only consider standard, $\ra$-discrete probability measure-preserving groupoids and standard probability spaces.
\end{description}

In this Section we will prove (Theorem \ref{theoremsecondimportant}) that, under weak dynamical conditions, the  full group of a probability measure-preserving groupoid, as a metric group, is enough to determine soficity. This answer, in this case, a question posed by Conley, Kechris and Tucker-Drob in \cite{MR3035288}.

If $(X,\mu)$ is a probability space, seen as a unit groupoid, then the measured semigroup $\MEAS(X,\mu)$ coincides with what is usually called its \emph{measure algebra}, that is, the $\sigma$-complete Boolean algebra of Borel subsets of $X$ modulo null sets, endowed with its given probability measure. We will follow the usual convention and denote it by $\MAlg(X,\mu)$ (or simply $\MAlg(X)$) instead, in order to decrease the risk of confusion later on.\nomenclature[Z]{$\MAlg(X),\MAlg(X,\mu)$}{Measure algebra of $(X,\mu)$}\index{Measure algebra} Given a probability measure-preserving groupoid $(\G,\mu)$, the set of idempotents of $\MEAS(\G,\mu)$ can be identified with $\MAlg(\G[0],\mu)$.

Again, for each $n\in\mathbb{N}$, fix $Y_n$ a set with $n$ elements and consider the full equivalence relation $Y_n^2$, whose full group $[Y_n^2]$ can be identified as the permutation group $\mathfrak{S}_n$ on $n$ elements, as in Example \ref{examplefullgroupofequivalencerelation}, with the normalized Hamming distance
\[d_\#(f,g)=\#\left\{i\in Y_n:f(i)\neq g(i)\right\}/n.\]\nomenclature[Z]{$\mathfrak{S}_n$}{Permutation group on $n$ elements}
A well-known theorem of Dye \cite{MR0158048} states that when $R$ is an aperiodic equivalence relation, the full group $[R]$ completely determines $R$. With this in mind, we prove that a probability measure-preserving groupoid $\G$ is sofic if and only if $[\G]$ embeds isometrically into $\prod_{\mathcal{U}}\mathfrak{S}_n$, as long as $\G$ is ``sufficiently dynamic''.

\begin{definition}\index{Metrically sofic group}
A group $G$ with a bi-invariant metric $d$ is \emph{metrically sofic} if it embeds isometrically into an ultraproduct $\prod_\mathcal{U}\mathfrak{S}_n$ of permutation groups with their normalized Hamming distances.
\end{definition}

\begin{example}
It is easy to come up with non-metrically sofic groups, even of diameter one. For example, let $G=\left\{1,i,-1,-i\right\}$ be the copy of the cyclic group of order $4$ inside $\mathbb{C}$, and endow it with the metric $d_G(g,h)=|g-h|/2$, which gives it diameter $1$. For every element $f$ of a permutation group $\mathfrak{S}_n$, $d_\#(f^2,1)\leq d(f,1)$, so the same is valid for elements of ultraproducts of permutation groups by \L o\'{s}' Theorem (\ref{lostheorem}). However, $d_G(i^2,1)>d(i,1)$, so $(G,d_G)$ is not metrically sofic.
\end{example}

We will need a few technical lemmas relating the full group $[\G]$, the measure algebra $\MAlg(\G[0])$ and the measured semigroup $\MEAS(\G)$. Before this, we do a quick detour back to general Boolean inverse monoid theory.

\begin{definition}\nomenclature[S]{$\supp(s)$}{Support of $s$}\index{Support}
If $S$ is a Boolean inverse monoid and $s\in S$, the \emph{support} of $s$ is $\supp(s)=s^*s\setminus\Fix(s)$.
\end{definition}

\begin{example}\label{examplesupportix}
If $S=\mathcal{I}(X)$, then $\supp(f)$ can be identified with its usual support, that is, $\left\{x\in\dom(f):f(x)\neq x\right\}$.
\end{example}

\begin{lemma}\label{lemmasupportssofic}
Let $S$ be a Boolean inverse monoid with a faithful normalized invariant mean $\mu$, and suppose $s,t\in S$, $ss^*=s^*s=tt^*=t^*t=1$. Then
\begin{enumerate}
\item[(a)] $d_\mu(s,1)=\mu(\supp(s))$;
\item[(b)] $d_\mu(s,t)=d_\mu(s^*t,1)$;
\item[(c)] $\supp(s)=\Fix(t)$ if and only if $d_\mu(s,t)=1$ and $\tr_\mu(s)+\tr_\mu(t)=1$;
\item[(d)] $\supp(s)\supp(t)=0$ if and only if $d_\mu(s,t)=d_\mu(1,s)+d_\mu(1,t)$. In this case, $\supp(s^*t)=\supp(s)\lor\supp(t)$;
\item[(e)] If $\supp(s)\supp(t)=0$ then $\supp(st)=\supp(s)\lor\supp(t)$.
\item[(f)] $\supp(sts^*)=s\supp(t)s^*$
\end{enumerate}
\end{lemma}
\begin{proof}
\begin{enumerate}
\item[(a)] This follows from the original definition of the uniform metric:
\[d_\mu(s,1)=\mu(s^*s\lor 1^*1\setminus\Fix(s^*1))=\mu(1\setminus\Fix(s))=1-\tr_\mu(s).\]
\item[(b)] Since $t^*t=tt^*=1$,
\[d_\mu(s,t)=d_\mu(st^*t,t)\leq d_\mu(st^*,1)=d_\mu(st^*,tt^*)\leq d_\mu(s,t).\]
\item[(c)] First suppose $d_\mu(s,t)=\tr_\mu(s)+\tr_\mu(t)=1$. Since $\Fix(s)\Fix(t)\leq\Fix(s^*t)$, then (a) and (b) imply $d_\mu(s,t)\leq 1-\mu(\Fix(s)\Fix(t))$, so since $\mu$ is faithful, we have $\Fix(s)\Fix(t)=0$. From the second equality, we conclude that $\mu(\Fix(s)\lor\Fix(t))=1$, and therefore $\Fix(t)=1\setminus\Fix(s)=\supp(s)$.

In the other direction, suppose $\supp(s)=\Fix(t)$. Then of course $\tr_\mu(s)+\tr_\mu(t)=1$, and for the other equality, we simply prove that $\Fix(s^*t)=0$. Indeed, if $e=e^2\leq s^*t$, then
\[\Fix(t)e=s^*t\Fix(t)e=s^*\Fix(t)e\]
which implies $\Fix(t)e\leq\Fix(s^*)=\Fix(s)=1\setminus\Fix(t)$. But also $\Fix(t)e\leq\Fix(t)$, so $\Fix(t)e=0$. Similarly, $\Fix(s)e=0$, and therefore $e=e(\Fix(s)\lor\Fix(t))=0$.
\item[(d)] Since $\Fix(s)\Fix(t)\leq\Fix(s^*t)$, then taking complements in $E(S)$ yields
\[\supp(s^*t)\leq\supp(s)\lor\supp(t)\ntag\label{equationlemmasupportssofic}\]
and then
\begin{align*}
d_\mu(s,t)&=\mu(\supp(s^*t))\leq\mu(\supp(s)\lor\supp(t))\\
&\leq \mu(\supp(s))+\mu(\supp(t))=d_\mu(s,1)+d_\mu(t,1).\ntag\label{lemmasupportssofic2}
\end{align*}
Note that the second inequality is an equality if and only if $\supp(s)\supp(t)=0$. This proves that if $d_\mu(s,1)+d_\mu(t,1)=d_\mu(s,t)$ then $\supp(s)\supp(t)=0$.

Conversely, if $\supp(s)\supp(t)=0$ then $\Fix(s)\lor\Fix(t)=1$. Since
\begin{align*}
\Fix(t)\Fix(s^*t)=\Fix(s^*t)\Fix(t)=s^*t\Fix(s^*t)\Fix(t)\\
=s^*t\Fix(t)\Fix(s^*t)=s^*\Fix(t)\Fix(s^*t),
\end{align*}
Then $\Fix(t)\Fix(s^*t)\leq\Fix(s^*)=\Fix(s)$, therefore
\[\Fix(s^*t)=\Fix(s)\Fix(s^*t)\lor\Fix(t)\Fix(s^*t)\subseteq\Fix(s)\]
and similarly $\Fix(s^*t)\leq\Fix(t)$. This proves $\Fix(s^*t)=\Fix(s)\Fix(t)$, so $\supp(s)\lor\supp(t)=\supp(s^*t)$, which is the opposite inclusion of Equation (\ref{equationlemmasupportssofic}), and so all inequalities in Equation (\ref{lemmasupportssofic2}) are in fact equalities.
\item[(e)] This follows from (d) and the fact that $\supp(s)=\supp(s^*)$, because $s^*s=ss^*=1$.
\item[(f)] $\supp(sts^*)=1\setminus sts^*=ss^*\setminus sts^*=s(1\setminus t)s^*=s\supp(t)s^*$.\qedhere
\end{enumerate}
\end{proof}

From items (c) and (d) of the lemma above, we immediately obtain:

\begin{corollary}\label{corollarysupportssofic}
Let $\theta:[\G]\to\prod_{\mathcal{U}}\mathfrak{S}_n$ be an isometric embedding and $\alpha,\beta\in[\G]$.
\begin{enumerate}
\item[(a)] $\supp(\alpha)=\Fix(\beta)$ if and only if $\supp(\theta(\alpha))=\Fix(\theta(\beta))$.
\item[(b)] $\supp(\alpha)\cap\supp(\beta)=\varnothing$ if and only if $\supp(\theta(\alpha))\land\supp(\theta(\beta))=0$.
\end{enumerate}
\end{corollary}

\begin{definition}\index{Groupoid!Aperiodic}
A probability measure-preserving groupoid $\G$ is \emph{aperiodic} if $\G_x$ is infinite for a.e.\ $x\in \G[0]$.
\end{definition}

\begin{example}
A standard probability measure-preserving equivalence relation $R$ on a standard probability space $(X,\mu)$ is aperiodic if and only if the equivalence class of a.e.\ point of $x\in X$ is infinite.
\end{example}

\begin{lemma}\label{lemmaaperiodicsupport}
Suppose $(\G,\mu)$ is an aperiodic probability measure-preserving groupoid. Then for all $A\in\MAlg(\G[0])$, there exists $\alpha\in[\G]$ such that $\supp(\alpha)=A$.
\end{lemma}

Recall that we only consider $\ra$-discrete and standard groupoids.

\begin{proof}
We can decompose $\G$ as $\G=\H\sqcup\mathcal{K}$, where $\H$ and $\mathcal{K}$ are subgroupoids, $|\ra\left(\G_x\right)|=\infty$ for all $x\in\H^{(0)}$, and $\mathcal{K}^x_x$ is infinite for all $x\in\mathcal{K}^{(0)}$, and we may assume that both $\H[0]$ and $\mathcal{K}^{(0)}$ are non-null in $\G[0]$ (in either case the proof can be adapted in an obvious manner).

The equivalence relation $\mathcal{R}(\H)=(\ra,\so)(\H)$ on $\H^{(0)}$ is also aperiodic, so by \cite[Lemma 4.10]{MR2583950}, there exists $f\in[\mathcal{R}(\H)]$ such that $\supp(f)=A\cap\H[0]$, and by Theorem \ref{theoremquotientfrommeasuredgroupoidtoorbitrelation} there exists $\alpha_{\H}$ such that $(\ra,\so)(\alpha_{\H})=f$, so $A\cap\H[0]=\supp(\alpha_{\H})$.

Consider the subset of $\mathcal{K}$
\[Z=\left\{k\in \mathcal{K}:\so(k)=\ra(k)\right\}\setminus A.\]
The restriction of the source map to $Z$ is countable-to-one and its image contains $A\cap\mathcal{K}[0]$, so by Corollary \ref{corollarysurjectiveborelmapshavesections} implies that there exits $\beta\subseteq Z$ such that $\so(\beta)=A$. Then $\alpha_{\mathcal{K}}=\beta\cup(\mathcal{K}\setminus A)$ belongs to the full group $[\mathcal{K}]$, and $\supp(\alpha_\mathcal{K})=A\cap\mathcal{K}^{(0)}$

To finish, let $\alpha=\alpha_{\H}\cup\alpha_{\mathcal{K}}$, so $\alpha\in[\G]$ and $\supp(\alpha)=A$.\qedhere
\end{proof}

\begin{lemma}\label{lemmasupportspreservedunderembeddings}
Suppose that $(\G,\mu)$ is an aperiodic probability measure-preserving groupoid, $\theta:[\G]\to\prod_{\mathcal{U}}\mathfrak{S}_n$ is an isometric embedding, and let $\alpha,\beta\in[\G]$. Then $\supp(\alpha)=\supp(\beta)$ if and only if $\supp(\theta(\alpha))=\supp(\theta(\beta))$.
\end{lemma}
\begin{proof}
Assume $\supp(\alpha)=\supp(\beta)$. By Lemma \ref{lemmaaperiodicsupport}, choose $\gamma\in[\G]$ such that $\supp(\gamma)=X\setminus\supp(\alpha)$, i.e., $\Fix(\gamma)=\supp(\alpha)$. Applying Corollary \ref{corollarysupportssofic}(a) twice,
\begin{align*}
\supp(\alpha)=\supp(\beta)&\iff\Fix(\gamma)=\supp(\theta(\beta))\iff\Fix(\theta(\gamma))=\supp(\theta(\beta))\\
&\iff\supp(\theta(\alpha))=\supp(\theta(\beta)).\qedhere
\end{align*}
\end{proof}

The next lemma will allow us to extend elements of $\MEAS(\G)$ to elements of $[\G]$.

\begin{lemma}\label{lemmaextensionfullsemigrouptofullgroup}
If $\gamma\in\MEAS(\G)$, then there exists $\widetilde{\gamma}\in[\G]$ with $\gamma\subseteq\widetilde{\gamma}$.
\end{lemma}
\begin{proof}
Let $\gamma_0=\gamma$, and for all $n\geq 1$, set
\[\gamma_n=\left\{g\in \gamma^{-n}:\so(g)\not\in \so(\gamma)\text{ and }\ra(g)\not\in \ra(\gamma)\right\}\]
Using the fact that $\gamma$ is a bisection and by construction of the $\gamma_n$, we have
\[\so(\gamma_n)\cap\so(\gamma_m)=\ra(\gamma_n)\cap\ra(\gamma_m)=\varnothing,\qquad\text{whenever }n\neq m.\]

The sets $\bigcup_n\so(\gamma_n)$ and $\bigcup_n\ra(\gamma_n)$ are both contained $\so(\gamma)\cup\ra(\gamma)$, so let us prove they are conull in there. Set
\[B=\ra(\gamma)\setminus \bigcup_n\so(\gamma_n)=\bigcap_n\ra(\gamma^n)\setminus\so(\gamma),\]
and for all $m\geq 0$, set
\[B_m=\ra(\gamma^{-m}B)=\bigcap_n\ra(\gamma^n)\cap\so(\gamma^m)\setminus\so(\gamma^{m+1})\]
These sets are pairwise disjoint, and since $\mu$ is invariant then they have the same measure, because
\[B_{m+1}=\ra(\gamma^{-1}B_m),\quad\text{and}\quad B_m=\ra(\gamma B_{m+1}),\quad\text{for all }m.\]
But $\mu$ is finite, therefore $\mu(B)=0$, and $\bigcup_n\so(\gamma_n)$ is conull in $\so(\gamma)\cup\ra(\gamma)$. Similarly, $\bigcup_n\ra(\gamma_n)$ is conull in $\so(\gamma)\cup\ra(\gamma)$. 

Therefore $\widetilde{\gamma}=\bigcup_n\gamma_n\cup(\G[0]\setminus(\so(\gamma)\cup\ra(\gamma))$ has the desired properties.\qedhere
\end{proof}

\begin{example}
If $R$ is a $\ra$-discrete probability-measure preserving Borel equivalence relation on a standard probability space $(X,\mu)$, then the previous lemma implies that every partial Borel automorphism $f:A\to B$ ($A,B\subseteq X$ Borel) with $\operatorname{graph}(f)\subseteq R$ can be extended (up to a null subset of $A$) to an a.e.-defined Borel automorphism $\widetilde{f}:X\to X$ with $\operatorname{graph}(f)\subseteq R$.
\end{example}

For the next lemma, again consider finite sets $Y_n$ with $\#Y_n=n$. We need to check that the group of invertible elements of $\prod_\mathcal{U}\MEAS(Y_n^2)$ (that is, those elements $\sigma$ for which $\sigma\sigma^*=\sigma^*\sigma=1$) coincides with $\prod_\mathcal{U}\mathfrak{S}_n$.

Recall that an element of the ultraproduct of metric structures $\prod_\mathcal{U}M_n$ is denoted as $(x_n)_\mathcal{U}$, where $x_n\in M_n$ for all $n$.

\begin{lemma}\label{lemmainvertibleinultraproduct}
If $\sigma=(\sigma_n)\in\prod_\mathcal{U}\MEAS(Y_n^2)$ is invertible, in the sense that $\sigma\sigma^*=(\id_{Y_n})_{\mathcal{U}}$, then there are $\widetilde{\sigma}_n\in\mathfrak{S}_n$ such that $\sigma=(\widetilde{\sigma}_n)_\mathcal{U}$.
\end{lemma}
\begin{proof}
Since $d_\#(\id_{Y_n},\sigma_n\sigma_n^*)=1-\mu_\#(\sigma_n^*\sigma_n)$ converges to $0$ along $\mathcal{U}$, then $\mu_\#(\sigma_n^*\sigma_n)$ converges to $1$. Simply extend each $\sigma_n$ arbitrarily to an element $\widetilde{\sigma}_n\in\mathfrak{S}_n$, so that
\[d_\#(\widetilde{\sigma}_n,\sigma)=\mu(\widetilde{\sigma}_n^*\widetilde{\sigma}_n\setminus\sigma_n^*\sigma_n)=1-\mu_\#(\sigma_n^*\sigma_n)\]
which, again, converges to $0$ along $\mathcal{U}$, and therefore $(\sigma_n)_\mathcal{U}=(\widetilde{\sigma}_n)_\mathcal{U}$.\qedhere
\end{proof}

This way, the group of invertible elements of $\prod_\mathcal{U}\MEAS(Y_n^2)$ is naturally identified with $\prod_\mathcal{U}\mathfrak{S}_n$, whereas the idempotent lattice $E(\prod_\mathcal{U}\MEAS(Y_n^2))$ is identified with $\prod_\mathcal{U}\MAlg(Y_n)$ (because $E(\prod_\mathcal{U}\MEAS(Y_n^2))$ is simply the image of the idempotent lattice of $\prod_{n\in\mathbb{N}}\MEAS(Y_n^2)$ under the canonical quotient map; see Subsection \ref{subsectionultraproductsofbooleaninversemonoids}).

\begin{definition}\index{Covariant}\label{definitioncovariant}
If $G_1$ and $G_2$ are groups acting on sets $X_1$ and $X_2$, respectively, $\theta:G_1\to G_2$ is a morphism and $\phi:X_1\to X_2$ is a function, we say that the pair $(\theta,\phi)$ is \emph{covariant} if $\phi(g x)=\theta(g)\phi(x)$ for all $g\in G_1$ and $x\in X_1$. If $G_1=G_2=G$ and $\theta=\id_G$, we simply say that  $\phi$ is covariant.
\end{definition}

\begin{remark}
Covariant pairs will be used also in Section \ref{sectionsoficactions} and in Chapter \ref{chapterpartialactionsofinversesemigroups} (see Definition \ref{definitioncovariantrepresentation}). In fact, covariant maps are a useful tool when constructing morphisms from structures which are ``built from'' both $G$ and a $G$-space $X$: In fact, the proofs of Theorems \ref{theoremfirstimportant} and \ref{theoremuniversalpropertysteinberg} follow similar general lines in different settings (however the fine details in each case differ substantially).
\end{remark}

\begin{definition}\label{definitionactionofinvertibleones}
Let $S$ be a Boolean inverse monoid and consider $S^\times$ the group of $g\in S$ such that $gg^*=g^*g=1$. We define an action of $S^\times$ on $E(S)$ as $g\cdot e=geg^*$. Note that if $\mu$ is an invariant mean on $S$ then for all $g\in S^\times$ and $e\in E(S)$,
\[\mu(e)=\mu(eg^*ge)=\mu((ge)^*(ge))=\mu((ge)(ge)^*)=\mu(geg^*)\]
so this action is isometric and trace-preserving. In particular,
\begin{itemize}
\item If $S=\MEAS(\G)$ for a probability measure-preserving groupoid $\G$, then $[\G]$ acts on $\MAlg(\G[0])$ via
\[\gamma\cdot A=\ra(\gamma A),\qquad\text{for all }\gamma\in[\G]\text{ and }A\in\MAlg(\G[0]).\]
\item If $S=\prod_\mathcal{U}\MEAS(Y_n^2)$ (where $Y_n$ are finite sets with $n$ elements), then $\prod_\mathcal{U}\mathfrak{S}_n$ acts on $\prod_\mathcal{U}\MAlg(Y_n)$ via
\[\left(\sigma_n\right)_\mathcal{U}\cdot\left(A_n\right)_\mathcal{U}=\left(\sigma_n(A_n)\right)_\mathcal{U}.\]
\end{itemize}
\end{definition}

\begin{theorem}\label{theoremfirstimportant}
A standard, $\ra$-discrete, aperiodic probability measure-preserving groupoid $(\G,\mu)$ is sofic if and only if its full group $[\G]$ is metrically sofic. Moreover, every isometric group embedding of $[\G]$ into an ultraproduct $\prod_\mathcal{U}\mathfrak{S}_n$ extends uniquely to a sofic embedding of $\G$.
\end{theorem}

Before starting, let us give an idea of the extension part of the proof (which is the hardest): Using the previous lemmas, every element of $\MEAS(\G)$ can be written as $\alpha A$, where $\alpha\in[\G]$ and $A\in\MAlg(\G[0])$, so if we have two semigroup embeddings $\theta:[\G]\to\prod_\mathcal{U}\mathfrak{S}_n$ and $\phi:\MAlg(\G[0])\to\prod_\mathcal{U}\MAlg(Y_n)$, we may define $\Theta(\alpha A)=\theta(\alpha)\phi(A)$, which a priori might depend on the choice of $\alpha$ and $A$. To guarantee that $\Theta$ is a morphism, we $(\theta,\phi)$ to be covariant. To obtain this, we will define $\phi$ \emph{from} $\theta$.

\begin{proof}[Proof of \ref{theoremfirstimportant}]
If $\G$ is sofic, then any sofic embedding $\Theta:\MEAS(\G)\to\prod_\mathcal{U}\MEAS(Y_n^2)$ restricts to an isometric group embedding of $[\G]$, which is therefore metrically sofic.

Let us prove the uniqueness first: If $\Theta:\MEAS(\G)\to\prod_\mathcal{U}\MEAS(Y_n^2)$ is an isometric embedding, then for every $\gamma\in\MEAS(\G)$ choose, by Lemmas \ref{lemmaaperiodicsupport} and \ref{lemmaextensionfullsemigrouptofullgroup}, $\widetilde{\gamma}$ and $\beta\in [\G]$ with $\supp(\beta)=\so(\gamma)$ and $\gamma\subseteq\widetilde{\gamma}$. By Theorem \ref{theoremtracepreservingisisometric}, $\Theta$ is a unital morphism of Boolean inverse semigroups and in particular it preserves supports, so
\[\Theta(\gamma)=\Theta(\widetilde{\gamma})\supp\Theta(\beta),\]
hence $\Theta$ is uniquely determined by its restriction to $[\G]$.

Now suppose $\theta:[\G]\to\prod_\mathcal{U}\mathfrak{S}_n$ is an isometric embedding, and let us use the ideas above to extend it to $\MEAS(\G)$.

We define $\phi:\MAlg(\G[0])\to\prod_\mathcal{U}\MAlg(Y_n)$ as follows: given $A\in\MAlg(\G[0])$, choose, by Lemma \ref{lemmaaperiodicsupport}, $\alpha\in[\G]$ such that $\supp(\alpha)=A$, and define $\phi(A)=\supp(\theta(A))$. By Lemma \ref{lemmasupportspreservedunderembeddings}, $\phi(A)$ does not depend on the choice of $\alpha$. We need several steps to finish this proof, namely,

%%\begin{itemize}
%%\item[1.]{\bf$\phi$ is well-defined, i.e., $\phi(A)$ does not depend on the choice of $\alpha$ with $\supp(\alpha)=A$:}
%%\end{itemize}
%%
%%Suppose $\alpha,\alpha'\in[\G]$ satisfy $\supp(\alpha)=\supp(\alpha')=A$. Consider any $\beta\in[\G]$ with $\supp(\beta)=\G[0]\setminus A$. By Lemma \ref{corollarysupportssofic}(a),
%%\[\supp(\theta(\alpha))=\Fix(\theta(\beta))=\supp(\theta(\alpha)).\]
%
%\item\label{itempreservesdisjointness} $\phi$ preserves disjointness, by 
%
%Suppose $A\cap B=\varnothing$. Choose $\alpha,\beta,\gamma\in[G]$ with $\supp(\alpha)=A$, $\supp(\beta)=B$, $\supp(\gamma)=G^{(0)}\setminus(A\cup B)$. Here we consider representatives $\theta(\alpha)=(a_n)_\mathcal{U}$, $\theta(\beta)=(b_n)_\mathcal{U}$, $\theta(\gamma)=(c_n)_\mathcal{U}$. By the previous Lemma again, we can approximate, for $\mathcal{U}$-a.e.\ $n$, $\mu_\#(\supp(c_nb_n)\triangle\Fix(a_n))\sim 0$, so $b_n=c_n^{-1}$ in $\supp(a_n)\cap\supp(b_n)$ up to a set of measure $\sim 0$, and similarly $a_n=c_n^{-1}$ in $\supp(a_n)\cap\supp(b_n)$ up to a set of measure $\sim 0$. Thus
%\begin{align*}
%d_\#(a_n,b_n)&\sim\mu_\#(\supp(a_n)\triangle\supp(b_n))\\
%&=\mu_\#(\supp(a_n))+\mu_\#(\supp(b_n))-2\mu_\#(\supp(a_n)\cap\supp(b_n))\\
%&=d_\#(a_n,1)+d_\#(b_n,1)-2\mu_\#(\supp(a_n)\cap\supp(b_n).
%\end{align*}
%Taking the limit over $\mathcal{U}$, we have
%\[d_\#(\theta(\alpha),\theta(\beta))=d_\#(\theta(\alpha),1)+d_\#(\theta(\beta),1)-2\mu_{\#}(\phi(A)\cap\phi(B)).\]
%On the other hand $d(\alpha,\beta)=d(\alpha,1)+d(\beta,1)$, thus $\mu_\#(\phi(A)\cap\phi(B))=0$, which means that $\phi(A)\cap\phi(B)=\varnothing$.
%
\begin{enumerate}
\item{\bf$\phi$ preserves meets:}

Let $A,B\in\MAlg(\G[0])$, and consider $\alpha,\beta,\gamma\in[\G]$ with
\[\supp(\gamma)=A\cap B,\quad \supp(\alpha)=A\setminus B,\quad\text{and}\quad\supp(\beta)=B\setminus A.\]
By \ref{corollarysupportssofic}(b) and \ref{lemmasupportssofic}(e), the supports of $\theta(\gamma)$ and $\theta(\alpha)$ are disjoint and
\[\supp(\theta(\gamma\alpha))=\supp(\theta(\gamma)\theta(\alpha))=\supp(\theta(\gamma))\lor\supp(\theta(\alpha)),\]
and similarly for $\gamma$ and $\beta$. Lemma \ref{lemmasupportssofic}(e) also implies that
\[\supp(\gamma\alpha)=A\quad\text{and}\quad\supp(\gamma\beta)=B,\]
so again by \ref{corollarysupportssofic}(b), and using distributivity of meets and joins,
\begin{align*}
\phi(A)\land\phi(B)&=\supp(\theta(\gamma\alpha))\land\supp(\theta(\gamma\beta))\\
&=\left(\supp(\theta(\gamma))\lor\supp(\theta(\alpha))\right)\land\left(\supp(\theta(\gamma))\lor\supp(\theta(\beta))\right)\\
%&=\supp(\theta(\gamma))\lor\left(\supp(\theta(\gamma))\land\supp(\theta(\beta))\right)\lor\left(\supp(\theta(\alpha))\land\supp(\theta(\gamma))\right)\lor\left(\supp(\theta(\alpha))\land\supp(\theta(\beta))\right)\\
&=\supp(\theta(\gamma))=\phi(A\cap B).
\end{align*}
%

\item{\bf If $\alpha\in{[\G]}$, then $\phi(\Fix(\alpha))=\Fix(\theta(\alpha))$:}

Choose again $\beta\in[\G]$ with $\supp(\beta)=\Fix(\alpha)$, so by \ref{corollarysupportssofic}(a),
\[\phi(\Fix(\alpha))=\phi(\supp(\beta))=\supp(\theta(\beta))=\Fix(\theta(\alpha)).\]
%

\item{\bf$\phi$ preserves the respective invariant means:}

Since $\theta$ is isometric then it is trace-preserving, so $\phi$ preserves invariant means by item 2.

\item{\bf$(\theta,\phi)$ is covariant:}

Let $A\in\MAlg(\G[0])$ and $\alpha\in[\G]$. Take $\beta\in[\G]$ with $\supp\beta=A$. Using \ref{lemmasupportssofic}(f),
\begin{align*}
\phi(\alpha\cdot A)&=\phi(\supp(\alpha\beta\alpha^*))=\supp(\theta(\alpha\beta\alpha^*))=\theta(\alpha)\cdot\supp(\theta(\beta))=\theta(\alpha)\cdot\phi(A).
\end{align*}

Recall that both $[\G]$ and $\MAlg(\G[0])$ are subsemigroups of $\MEAS(\G)$, and similarly for $\prod_\mathcal{U}\mathfrak{S}_n$ and $\prod_\mathcal{U}\MAlg(Y_n)$ inside $\prod_\mathcal{U}\MEAS(Y_n^2)$.

\item{\bf  If $\alpha,\beta\in[\G]$, $A\in\MAlg(\G[0])$ satisfy $\alpha A=\beta A$, then $\theta(\alpha)\phi(A)=\theta(\beta)\phi(A)$:}

%Since $\alpha^*\alpha=\beta^*\beta=\G[0]$ and $A$ and $B$ are idempotents,
%\[A=A^*\alpha^*\alpha A=(\alpha A)^*(\alpha A)=(\beta B)^*(\beta B)=B\]
%thus
%\[\alpha^*\beta A=\alpha^* \alpha A=A,\]
%so $A\subseteq\Fix(\alpha^*\beta)$.
Items 1.\ and 2.\ above imply
\[\phi(A)\subseteq\Fix(\theta(\alpha^*\beta)),\]
which in turn yields
\[\theta(\beta)\phi(B)=\theta(\alpha)\theta(\alpha^*\beta)(\phi(A)=\theta(\alpha)\phi(A).\]
\end{enumerate}

Now we define $\Theta:\MEAS(\G)\to\prod_\mathcal{U}\MEAS(Y_n^2)$ as $\Theta(\alpha)=\theta(\widetilde{\alpha})\phi(\alpha^*\alpha)$, where $\widetilde{\alpha}\in[\G]$ is chosen, by Lemma \ref{lemmaextensionfullsemigrouptofullgroup}, such that $\alpha\subseteq\widetilde{\alpha}$. Item 5.\ above implies that $\Theta(\alpha)$ does not depend on the choice of $\widetilde{\alpha}$.
%For every $\alpha\in\MEAS(\G)$, choose $\widetilde{\alpha}\in[\G]$ with $\alpha\subseteq\widetilde{\alpha}$ and set$\Theta(\alpha)=\theta(\widetilde{\alpha})\phi(\alpha^*\alpha)$. In order to prove that $\Theta$ does not depend on the choice of $\widetilde{\alpha}$, suppose $\alpha\subseteq\widetilde{\beta}\in[\G]$. This means that
%\[\widetilde{\alpha}\alpha^*\alpha=\alpha=\widetilde{\beta}\alpha^*\alpha,\qquad\text{ so }\alpha^*\alpha\subseteq\Fix(\widetilde{\alpha}^*\widetilde{\beta}).\]
%By items 2.\ and 3.\ above, we conclude that
%\[\phi(\alpha^*\alpha)\subseteq\Fix(\theta(\widetilde{\alpha}^*\widetilde{\beta})),\]
%and this implies
%\[\theta(\widetilde{\beta})\phi(\alpha^*\alpha)=\theta(\widetilde{\alpha})\theta(\widetilde{\alpha^*}\widetilde{\beta})\phi(\alpha^*\alpha)=\theta(\widetilde{\alpha})\phi(\alpha^*\alpha).\]

Note that $\Theta$ extends both $\theta$ and $\phi$. Let us show that $\Theta$ is a sofic embedding.

Suppose $\alpha\subseteq\widetilde{\alpha}$, $\beta\subseteq\widetilde{\beta}$, where $\alpha,\beta\in\MEAS(\G)$ and $\widetilde{\alpha},\widetilde{\beta}\in[\G]$. Then
\[\alpha\beta=\widetilde{\alpha}\widetilde{\beta}\left(\beta^*\beta\cap(\beta^*\cdot\alpha^*\alpha)\right).\]
In fact, this formula holds independently of $\G$, so it will hold with $Y_n^2$ in place of $\G$ and on $\prod_\mathcal{U}\MEAS(Y_n^2)$ by \L o\'{s}' Theorem (\ref{lostheorem}). This is enough to conclude that $\Theta(\alpha\beta)=\Theta(\alpha)\Theta(\beta)$, since both $(\theta,\phi)$ is a covariant pair of semigroup morphisms.%Since a similar formula holds on ultraproducts and $(\theta,\phi)$ is covariant, we obtain $\Theta(\alpha\beta)=\Theta(\alpha)\Theta(\beta)$.

It remains only to check that $\Theta$ is trace-preserving. Let $\alpha\in\MEAS(\G)$. If we show that $\Fix(\Theta(\alpha))=\phi(\Fix\alpha)$ we are done because $\phi$ preserves invariant means.

Given $\widetilde{\alpha}\in[\G]$ with $\alpha\subseteq\widetilde{\alpha}$, let $A=\Fix(\alpha)$ and $B=\Fix(\widetilde{\alpha})\setminus A$. Note that $A=\Fix(\widetilde{\alpha})\cap\alpha^*\alpha$, so
\[\phi(A)=\Fix\left(\theta(\widetilde{\alpha})\right)\land (\phi(\alpha^*\alpha)).\]
Since $\Theta(\alpha)=\theta(\widetilde{\alpha})\phi(\alpha^*\alpha)\leq\theta(\widetilde{\alpha})$, then $\Theta(\alpha)^*\Theta(\alpha)=\phi(\alpha^*\alpha)$, so
\[\Fix(\Theta(\alpha))=\Fix(\theta(\widetilde{\alpha}))\land (\phi(\alpha^*\alpha))=\phi(A)=\phi(\Fix(\alpha)),\]
and we conclude that $\Theta$ is a sofic embedding of $\G$.\qedhere
\end{proof}

We will now extend this result allowing $\G$ to have periodic points, however we still need that $\G_x$ is not a singleton for a.e. $x\in\G[0]$. Set
\[\operatorname{Per}_{\geq 2}(\G)=\left\{x\in \G[0]:2\leq\#(\G_x)<\infty\right\}.\]

\begin{lemma}
There exists $\alpha\in[\G]$ with $\supp(\alpha)=\operatorname{Per}_{\geq 2}(\G)$.
\end{lemma}
\begin{proof}
Let $\mathcal{P}=\operatorname{Per}_{\geq 2}(\G)$. Substituting $\G$ by the subgroupoid $\mathcal{P}\G\mathcal{P}$ if necessary, we may assume $\mathcal{P}=\G[0]$. As in the proof of Lemma \ref{lemmaaperiodicsupport}, decompose $\G$ as $\G=\H\sqcup \mathcal{K}$, where $\H$ and $\mathcal{K}$ are subgroupoids of $\G$, $\H_x$ is a  is principal, and $\#(\mathcal{K}^x_x)\geq 2$ for a.e.\ $x\in\mathcal{K}^{(0)}$.

The last part of Corollary \ref{corollarysurjectiveborelmapshavesections} applied to the restriction of the source map
\[\so:\left\{k\in\mathcal{K}:\so(k)=\ra(k)\right\}\setminus\mathcal{K}^{(0)}\to\mathcal{K}^{(0)}\]
provides $\beta\in\MEAS(\mathcal{K})$ with $\supp(\beta)=\mathcal{K}^{(0)}$.

Since $\H$ is principal it is (isomorphic to) an equivalence relation on $\H[0]$. For each $n\geq 2$, let $\H_n=\left\{h\in\H:\#(\H_{\so(h)})=n\right\}$. Using the selection theorem for periodic relations \cite[Theorem 12.16]{MR1321597}, decompose $\H_n^{(0)}$ into a Borel partition $X_1,\ldots,X_n$ such that each $X_i$ contains exactly one representative of each $\H$-class. The Borel subset
\[\gamma_n=(X_1\times X_n)\cup \bigcup_{i=1}^{n-1}(X_{i+1}\times X_i)\]
is a bisection of $\H$. (Seen as a function, it maps each $x\in X_i$ to the unique element of $X_{i+1}$ (or $X_1$ if $i=n$) which is equivalent to $x$.) Then $\supp(\gamma_n)=\H[0]_n$.

Therefore $\alpha=\beta\cup\bigcup_n\gamma_n\in\MEAS(\G)$, and $\supp(\alpha)=\operatorname{Per}_{\geq 2}(\G)$.\qedhere
\end{proof}

\begin{theorem}\label{theoremsecondimportant}
Suppose $(\G,\mu)$ is a (standard, $\ra$-discrete) probability measure-preserving groupoid such that for all $x\in\G[0]$, $\#\G_x\geq 2$. Then $(\G,\mu)$ is sofic if and only if $[\G]$ is metrically sofic.
\end{theorem}
\begin{proof}
Let $P=\operatorname{Per}_{\geq 2}(\G)$ and $\operatorname{Aper}=\G[0]\setminus P$, and consider the subgroupoid $\H=\G\operatorname{Aper}$ of $\G$. We have already seen that the subgroupoid $\G P$ is sofic, since it is periodic, and since $\G$ is a convex combination of $\G P$ and $\H$, it suffices to show that $[\H,\mu_{\H}]$ is metrically sofic. Fix any $\rho\in[\G]$ with $\supp\rho=P$. Let $\theta:[\G]\to\prod_{\mathcal{U}}\mathfrak{S}_n$ be an isometric embedding.

Consider the embedding $\iota:[\H,\mu_{\H}]\to[\G,\mu]$, $\alpha\mapsto \widetilde{\alpha}=\alpha\cup P$. Then for all $\alpha\in[\H,\mu_{\H}]$,
\[\tr_\mu(\iota(\alpha))=\mu(\operatorname{Aper})\tr_{\H}(\alpha)\]

By \ref{corollarysupportssofic}(a), $\supp(\theta(\iota(\alpha)))\subseteq\Fix(\theta(\rho))$ whenever $\alpha\in[\H,\mu_{\H}]$.

As in the proof of Theorem \ref{theorempermanencesofic}(d), we can ``restrict'' $\theta(\iota(\alpha))$ to $\Fix(\theta(\rho))$ and obtain a new semigroup embedding $\eta:[\H]\to\prod_{\mathcal{U}}\mathfrak{S}_{m_n}$ (where $m_n\leq n$), in such a way that for all $\alpha\in[\H,\mu_{\H}]$,
\begin{align*}
\tr(\eta(\alpha))&=\tr(\theta(\rho))^{-1}\tr(\theta(\iota(\alpha)))=\tr_\mu(\rho)^{-1}\tr_\mu(\iota(\alpha))\\
&=\mu(\operatorname{Aper})^{-1}\mu(\operatorname{Aper})\tr_{\H}(\alpha)\\
&=\tr_{\H}(\alpha)\mu(\theta(\rho),1)^{-1}=d(\rho,1)^{-1}=\mu(\operatorname{Aper})^{-1}
\end{align*}
so $\eta$ is in fact a sofic embedding of $\H$.\qedhere
\end{proof}